% ============================================================================
%  电源线侧信道采集设备 — 设计提案
%  Power-Line Side-Channel Acquisition Device — Design Proposal
%
%  目标受众: U Michigan Yan Long Lab
%  文档类型: 工程设计提案(中文)
%  编译命令: xelatex -shell-escape proposal.tex (跑两遍以生成 TOC)
% ============================================================================
\documentclass[a4paper,11pt,oneside,UTF8,scheme=plain]{ctexart}

% ---------- 几何与版式 ----------
\usepackage[margin=2.5cm,top=2.5cm,bottom=2.5cm,headheight=14pt]{geometry}
\usepackage{setspace}
\onehalfspacing
\setlength{\parskip}{0.4em}

% ---------- 字体配置 ----------
% 中文用思源宋体(标题用黑体),拉丁字符在大字号下也用思源宋体覆盖避免 fallback 缺失
\setCJKmainfont[BoldFont={Source Han Serif CN Bold}]{Source Han Serif CN}
\setCJKsansfont[BoldFont={Source Han Sans CN Bold}]{Source Han Sans CN}
\setCJKmonofont{Source Han Sans CN}

\setmainfont[BoldFont={Source Han Serif CN Bold}]{Source Han Serif CN}
\setsansfont[BoldFont={Source Han Sans CN Bold}]{Source Han Sans CN}
\setmonofont[Scale=0.88]{JetBrains Maple Mono}

% ---------- 数学 ----------
\usepackage{amsmath,amssymb,amsthm}
\usepackage{unicode-math}
\setmathfont{Latin Modern Math}

% ---------- 表格 ----------
\usepackage{booktabs}    % \toprule \midrule \bottomrule
\usepackage{array}
\usepackage{tabularx}
\usepackage{longtable}
\usepackage{makecell}
\usepackage{multirow}
\renewcommand{\arraystretch}{1.20}

% ---------- 颜色(严肃学术: 仅保留黑/灰 + 极淡蓝灰,无强色) ----------
\usepackage[table]{xcolor}
% 统一一种深灰色作为标题/规则线颜色; 表头用极浅灰
\definecolor{rulegray}{HTML}{2D2D2D}      % 深灰,代替原先的 accent 蓝
\definecolor{tableheader}{HTML}{F5F5F5}   % 极浅灰表头
\definecolor{linkgray}{HTML}{555555}      % 超链接深灰
\definecolor{codebg}{HTML}{FAFAFA}        % 代码块极浅灰背景

% ---------- 代码块(去除花哨颜色) ----------
\usepackage{listings}
\lstdefinestyle{codestyle}{
  backgroundcolor=\color{codebg},
  basicstyle=\ttfamily\footnotesize,
  numbers=left,
  numberstyle=\tiny\color{gray!60},
  numbersep=8pt,
  frame=single,
  framerule=0.4pt,
  rulecolor=\color{gray!40},
  showstringspaces=false,
  breaklines=true,
  tabsize=4,
  captionpos=b,
  keywordstyle=\bfseries,
  commentstyle=\color{gray!60}\itshape,
  stringstyle=\itshape,
  xleftmargin=1.5em,
  framexleftmargin=1.5em,
  framexrightmargin=0pt,
  framextopmargin=4pt,
  framexbottommargin=4pt,
}
\lstset{style=codestyle}
\lstdefinelanguage{verilog}{
  keywords={module,endmodule,input,output,wire,reg,assign,always,posedge,negedge,if,else,begin,end,case,endcase,parameter,localparam,generate,endgenerate,for,signed,unsigned},
  morecomment=[l]//,
  morecomment=[s]{/*}{*/},
  morestring=[b]",
  sensitive=true,
}
\lstdefinelanguage{tcl}{
  keywords={proc,set,if,else,foreach,for,while,return,puts,catch,error,source,open_project,open_bd_design,create_bd_port,connect_bd_net,launch_runs,wait_on_run,write_bitstream,read_xdc,get_bd_pins,get_bd_cells,get_bd_ports},
  morecomment=[l]\#,
  morestring=[b]",
  sensitive=true,
}

% ---------- 超链接(深灰,无明显颜色) ----------
\usepackage[colorlinks=true,linkcolor=linkgray,urlcolor=linkgray,citecolor=linkgray,unicode=true,bookmarksopen=true,bookmarksnumbered=true]{hyperref}
\usepackage{bookmark}

% ---------- 页眉页脚 ----------
\usepackage{fancyhdr}
\pagestyle{fancy}
\fancyhf{}
\fancyhead[L]{\small\nouppercase\leftmark}
\fancyhead[R]{\small 电源线侧信道采集设备}
\fancyfoot[C]{\thepage}
\renewcommand{\headrulewidth}{0.4pt}
\renewcommand{\footrulewidth}{0pt}

% ---------- 章节标题样式(去掉规则线,黑字加粗,标准学术报告) ----------
\usepackage{titlesec}
\titleformat{\section}
  {\Large\bfseries}
  {\thesection}{1em}{}
\titleformat{\subsection}
  {\large\bfseries}
  {\thesubsection}{1em}{}
\titleformat{\subsubsection}
  {\normalsize\bfseries}
  {\thesubsubsection}{1em}{}

% ---------- 浮动体 ----------
\usepackage{float}
\usepackage{graphicx}
\graphicspath{{figures/}{../figures/}}
\usepackage{caption}
\captionsetup[table]{labelfont={bf},textfont={normalfont},justification=centering,skip=4pt}
\captionsetup[figure]{labelfont={bf},textfont={normalfont},justification=centering,skip=4pt}

% ---------- 状态徽章命令(去除颜色,纯文字) ----------
\newcommand{\statusok}{\textbf{[已完成]}}
\newcommand{\statuspartial}{\textbf{[部分完成]}}
\newcommand{\statusplan}{\textbf{[计划中]}}

% ---------- 信息框(浅灰边框,无填充色) ----------
\usepackage[most]{tcolorbox}
\newtcolorbox{infobox}[1][]{
  colback=white,
  colframe=gray!50,
  fonttitle=\bfseries,
  arc=0mm,
  boxrule=0.5pt,
  left=4pt,right=4pt,top=4pt,bottom=4pt,
  #1
}
\newtcolorbox{warnbox}[1][]{
  colback=white,
  colframe=gray!70,
  fonttitle=\bfseries,
  arc=0mm,
  boxrule=0.6pt,
  left=4pt,right=4pt,top=4pt,bottom=4pt,
  #1
}

% ---------- 标题信息 ----------
\title{\vspace{-2em}\sffamily\bfseries 电源线侧信道采集设备\\[0.3em]\large 设计提案 v1.0}
\author{}
\date{}

% ============================================================================
\begin{document}
% ============================================================================

% ---------- 封面页 ----------
\thispagestyle{empty}
\vspace*{2cm}

\begin{center}
{\sffamily\bfseries\Huge 电源线侧信道采集设备}\\[0.6em]
{\sffamily\Large 设\ 计\ 提\ 案}\\[1.5em]
{\large Power-Line Side-Channel Acquisition Device}\\[0.3em]
{\normalsize Design Proposal v1.0}\\[2em]

\rule{0.6\textwidth}{0.4pt}\\[1em]

{\large\sffamily 将 EM Eye 攻击从空中辐射通道延伸至电源线传导通道}\\[3em]

\begin{tabularx}{0.85\textwidth}{rX}
\sffamily\bfseries 申请岗位 & RA, Yan Long Lab \\
\sffamily\bfseries 申请单位 & University of Michigan \\
\sffamily\bfseries 候选人  & \textit{(待填写)} \\
\sffamily\bfseries 文档日期 & 2026 年 5 月 14 日 \\
\sffamily\bfseries GitHub  & \url{https://github.com/stongry/powerline-emeye-proposal} \\
\sffamily\bfseries 版本   & v1.0 (LaTeX 终稿) \\
\end{tabularx}

\vfill

\begin{infobox}[title=工程基础]
Pre-Phase 0（2026-05-12 至 2026-05-14）已完成：5 份 Python 仿真、4 份 RTL testbench 全 PASS、Vivado synth+impl+bitstream（\texttt{xc7z010clg400-2}）真目标烧录、AD9361 RF 实测（RX\_LO 204 MHz, RSSI 93.75 dB, 16 KB IQ）。详见附录 B。
\end{infobox}

\vspace{1em}
\end{center}

\newpage

% ---------- 目录 ----------
\tableofcontents
\newpage

% ============================================================================
\section*{摘要}
\addcontentsline{toc}{section}{摘要}
% ============================================================================

\begin{infobox}[title=核心论点]
CMOS 图像传感器及 MIPI 总线在工作时产生的高频电磁辐射，\textbf{以共模电流的形式沿电源线传导}，在远离目标设备的电源插座、电源排插甚至建筑配电线网的任意一点都可被拾取，将 EM Eye 风险从"近场 5 米"扩展到"配电网任意位置数十米"。
\end{infobox}

\paragraph{关键技术决策}

\begin{table}[H]
\centering
\caption{六项核心决策一览}
\begin{tabularx}{\textwidth}{l>{\raggedright\arraybackslash}X}
\toprule
\rowcolor{tableheader} \textbf{维度} & \textbf{决策} \\
\midrule
形态     & 便携集成式（\(150 \times 80 \times 50\) mm），BYO LDSDR 7010 模块，Phase 2 演进插头式 \\
耦合     & 主路宽带 CT（Tekbox TBCP2-1000，1 MHz – 1 GHz），辅路 HV 电容（覆盖 \(>\) 500 MHz） \\
模拟前端 & 候选人 BYO 自研 ERA-4SM+ \(\times\) 2 级联 LNA（VNA 实测 +28 dB @ 100 MHz） \\
采集     & LDSDR 7010（AD9363 \(\to\) AD9361 解锁，70 MHz – 6 GHz，12-bit IQ \(\times\) 2 RX，30 MSPS） \\
传输     & 千兆以太网直连主机，32 Mbps 双通道压缩流（板上 FPGA \texttt{emeye\_accel} 加速） \\
主控     & 直接使用 LDSDR 板内 Zynq-7010 PS，无需外接 RPi CM4（简化系统 + 降本约 ¥1000） \\
\bottomrule
\end{tabularx}
\end{table}

\paragraph{核心差异化} 候选人已在 Pluto Zynq-7010 平台完成 OFDM + LDPC 全链路 HDL 收发机（BER = 0 板级验证 \statusok），并在 Pre-Phase 0 阶段完成 RTL 设计、单元测试、Vivado 实现、真硬件 bitstream 烧录与 RF 链路实测（详见附录 B）。

\paragraph{交付计划}（6 个月，1 RA 全职）

\begin{table}[H]
\centering
\caption{分阶段交付里程碑}
\begin{tabularx}{\textwidth}{lcX}
\toprule
\rowcolor{tableheader} \textbf{阶段} & \textbf{周期} & \textbf{关键交付} \\
\midrule
Phase 0     & 3 周    & LISN + RPi V1 复现 EM Eye + 电源线传导验证 \(\to\) \textbf{Go/No-Go 决策门} \\
Phase 1     & 3 个月  & 耦合网络 + 模拟前端 + 集成原型 + Phase 0 修订 \\
Phase 2     & 2 个月  & 板上 FPGA \texttt{emeye\_accel} 加速 + 多设备验证 + 终版报告 \\
\rowcolor{gray!10} Phase 3 & {\color{gray}不在 6 个月预算内} & {\color{gray}多频段相干融合 + HDMI TEMPEST 泛化} \\
\bottomrule
\end{tabularx}\\
\footnotesize\textit{所有 \statusok 项目均有 git commit / 文件 / log / SSH session 输出佐证。}
\end{table}

\paragraph{预算概要}单台原型 BOM \(\approx\) ¥4{,}350（BYO 资产省 ¥3{,}000+）· NRE \(\approx\) ¥5{,}300 · 总预算 \(\approx\) ¥31{,}500。

\newpage

% ============================================================================
\section*{Pre-Phase 0 已完成工作速览}
\addcontentsline{toc}{section}{Pre-Phase 0 已完成工作速览}
% ============================================================================

本提案不是纸面方案。表 \ref{tab:prephase0-summary} 列出截至 2026-05-14 已经独立完成、可复现的工程工作清单。完整记录（代码、log、bitstream、SSH session 输出）详见附录 B。

\begin{table}[H]
\centering
\caption{Pre-Phase 0 已完成工作（2026-05-12 至 2026-05-14）}
\label{tab:prephase0-summary}
\small
\begin{tabularx}{\textwidth}{lXl}
\toprule
\rowcolor{tableheader}
\textbf{类别} & \textbf{产物 \& 关键数字} & \textbf{状态} \\
\midrule
Python 仿真 & EM Eye + HDMI TEMPEST 双场景，SSIM 0.98 / 0.99，1.2 s 运行 & \statusok \\
Verilog RTL & 4 模块 + 4 testbench 全 PASS，1156 输出 0 丢样 & \statusok \\
Vivado 流程 & xc7z010clg400-2 真目标 synth + impl + bitstream（2.0 MB） & \statusok \\
真硬件烧录 & LDSDR rev2.1 fpga\_manager 烧 .bin 成功，FPGA Manager operating & \statusok \\
RF 链路实测 & RX\_LO 204 MHz，RSSI 93.75 dB，16 KB IQ 真实采集 & \statusok \\
BD 改造     & ad9361\_phy\_0 ADC 端口暴露 + GPIO\_I 反馈通路 & \statusok \\
Phase 2 草稿 & BD tap + AXI-DMA + PS TCP forwarder 1{,}096 行文档 \& C 代码 & \statusok \\
\bottomrule
\end{tabularx}\\
\footnotesize\textit{GitHub：\url{https://github.com/stongry/powerline-emeye-proposal}}
\end{table}

\newpage

% ============================================================================
\section{总体形态、尺寸与系统框图}
\label{sec:overview}
% ============================================================================

\subsection{形态决策：Route A vs Route B}
\label{sec:form-factor}

约束：设备最大尺寸 \(\le 150 \times 80 \times 50\) mm。本提案锁定 \textbf{Route A 便携集成式}为 Phase 1 主交付形态，Route B 插头适配式作为 Phase 2 演进方向。

\begin{table}[H]
\centering
\caption{Route A vs Route B 形态对比}
\label{tab:routea-vs-b}
\small
\begin{tabularx}{\textwidth}{lXX}
\toprule
\rowcolor{tableheader}
\textbf{维度} & \textbf{Route A 便携集成式（Phase 1）} & \textbf{Route B 插头适配式（Phase 2）} \\
\midrule
供电方式       & USB-C 5 V（移动电源 / 笔记本 / 充电器） & 内含 AC-DC，直插市电 \\
与市电关系      & \textbf{电气完全隔离}（仅 CT 磁耦合）   & 直接接触 220 V \\
触电风险       & 归零                                 & 高，需 IEC 61010 + EMC Class B 认证 \\
保护链复杂度    & TVS + DC 隔直 + LC HPF（约 ¥11）       & 全套 GDT/MOV/共模扼流/TVS（约 ¥80+） \\
BOM 净成本     & \(\sim\) ¥820                        & \(\sim\) ¥1{,}300+ \\
开发周期       & 3 周 Phase 0 + 12 周 Phase 1          & + 4 周安规认证准备 \\
\bottomrule
\end{tabularx}
\end{table}

\textbf{选定 Route A 理由}：(1) 安全性：整机与市电零电气连接，触电风险归零；(2) 快速交付：Phase 1 三个月可出原型；(3) 形态自由：USB 供电覆盖实验室与现场便携。

\subsection{尺寸预算}
\label{sec:dim-budget}

外壳目标 \(100 \times 60 \times 30\) mm（留 50\% 边距，极限可压到 \(80 \times 50 \times 25\) mm），内部子模块布局如表 \ref{tab:dim-budget}。

\begin{table}[H]
\centering
\caption{Route A 内部尺寸分配}
\label{tab:dim-budget}
\small
\begin{tabularx}{\textwidth}{lccX}
\toprule
\rowcolor{tableheader}
\textbf{模块} & \textbf{尺寸 (mm)} & \textbf{占比} & \textbf{备注} \\
\midrule
模拟前端板（CT 二次侧 → LNA 输出） & \(80 \times 60 \times 8\) & 35\% & 4 层 PCB，RF 路径 Rogers RO4350B，LNA 屏蔽腔 \\
LDSDR 7010 rev2.1 SOM            & \(90 \times 50 \times 12\) & 38\% & 候选人 BYO，标准 Pluto 衍生板形态 \\
USB-C + ESD + TPS7A47 LDO        & \(30 \times 20 \times 6\)  & 5\%  & 模拟域净化电源 \\
Tekbox TBCP2-1000 CT             & 外挂 OD 约 100 & — & RG-316 短同轴接入，\textbf{不入主壳} \\
Ethernet 输出 RJ45               & \(16 \times 13 \times 6\)  & 4\%  & 千兆 \(\to\) PC \\
散热与屏蔽间隙                    & — & 18\% & 顶/底盖空隙 + EMI 衬垫 \\
\bottomrule
\end{tabularx}
\end{table}

\subsection{系统总框图}
\label{sec:block-diagram}

完整系统由 7 个功能模块串联组成（图 \ref{fig:block-diagram} 占位，由后续 SVG 替换）：

\begin{figure}[H]
\centering
\framebox[0.95\textwidth]{
\begin{minipage}{0.92\textwidth}
\centering\ttfamily\small
DUT(220V) → Tekbox CT → TVS+隔直 → LC HPF (30 MHz) → ERA-4SM+ x2 LNA (+28 dB)\\
→ LDSDR 7010 (AD9363 70M-6G, 2RX, 千兆 Eth) → PS Linux → PC\\[0.4em]
\rmfamily 7 个模块: 耦合 · 保护 · 滤波 · 放大 · 采集 · 主控 · 传输
\end{minipage}
}
\caption{系统总框图（v1.0 占位，待 SVG 替换）}
\label{fig:block-diagram}
\end{figure}

7 个模块：(1) \textbf{耦合} Tekbox TBCP2-1000 商用宽带 CT，L+N 同向穿芯，DM 抑制 30–40 dB；(2) \textbf{保护} Bourns CDSOT23-T05LC 低 C TVS（\(C < 1\) pF）+ Murata NP0 1 nF DC 隔直；(3) \textbf{滤波} LC 4 阶 Butterworth HPF，\(f_c = 30\) MHz，通带损耗 \(< 1\) dB，50 Hz 衰减 \(> 100\) dB；(4) \textbf{放大} BYO ERA-4SM+ \(\times 2\)，+28 dB @ 100 MHz 实测 \statusok；(5) \textbf{采集} LDSDR 7010 rev2.1，70 MHz – 6 GHz，12-bit IQ，2 RX 并行；(6) \textbf{主控} LDSDR 板内 Zynq-7010 PS（不外接 CM4）；(7) \textbf{传输} 千兆以太网原始 IQ 直传，Phase 2 演进为板上加速 + 8-bit 解调流。

\subsection{Phase 2 向插头形态演进路线}
\label{sec:phase2-route-b}

Route B 为终态隐蔽攻击形态，Phase 2 集成阶段演进路径如下：

\begin{table}[H]
\centering
\caption{Route B 演进步骤与工时}
\label{tab:route-b-steps}
\small
\begin{tabularx}{\textwidth}{cXc}
\toprule
\rowcolor{tableheader}
\textbf{步骤} & \textbf{工作内容} & \textbf{工时} \\
\midrule
1 & 设计内部 AC-DC（Mean Well IRM-10-5），替代 USB 输入 & 2 周 \\
2 & 补全市电保护链（GDT + MOV + 共模扼流圈）           & 1 周 \\
3 & 自研小型化 CT（Fair-Rite \#43+\#61 复合磁芯），压到 OD \(\le 20\) mm，集成进壳内 & 3 周 \\
4 & 6 层 PCB 重新布板，IEC 61010 间距规则，加铜罩屏蔽 & 2 周 \\
5 & EMC Class B 预测试                                & 1 周 \\
\bottomrule
\end{tabularx}
\end{table}

Phase 2 总工时约 9 周，在 8 周（Week 17–24）预算内紧凑可行。

% ============================================================================
\section{电源线高频泄漏耦合与提取}
\label{sec:ch2}
% ============================================================================

\subsection{物理机制}
\label{sec:ch2-phy}

\subsubsection{CMOS + MIPI 总线产生共模电流的物理过程}

EM Eye 论文 §III：MIPI CSI-2 串行链路传输 sub-ns 边沿的高速数据，数据信号耦合到外围金属结构，连接电缆本身成为非预期发射天线（\textit{unintentional transmission antenna}）。物理机制三阶段：

\begin{enumerate}
\item \textbf{源}：MIPI CSI-2 D-PHY 以 byte clock（RPi V1 \(51\,\mathrm{MHz}\)）切换差分对，每次比特跳变 \(\mathrm{d}V/\mathrm{d}t \sim 1\,\mathrm{V/ns}\) \(\to\) 共模电感耦合产生共模电流分量。
\item \textbf{传播}：共模电流沿数据线 \(\to\) 主板地参考 \(\to\) PCB 寄生电容 \(\to\) 电源输入端 \(\to\) 稳压模块寄生回路 \(\to\) 外部电源线（L+N+PE）。
\item \textbf{泄漏频段}：byte clock 谐波富集，典型 100 MHz – 1 GHz 段，与论文 Table II 12 款 COTS 设备（155–1740 MHz）一致。
\end{enumerate}

\subsubsection{电源线作为“非预期传输天线”的延伸}

EM Eye 原论文使用近场磁探头（Beehive 100C）或 LPDA 远场天线在空气中拾取泄漏。共模电流同样沿电源线传导，依据如下：

\begin{itemize}
\item 电源线在 100 MHz – 1 GHz 段电气长度远大于 1/4 波长（0.5 m 线在 500 MHz 处 \(\approx 0.83\lambda\)），表现为分布参数传输线。
\item EMC Class B 电源滤波器典型工作到 30 MHz 即截止，对 100 MHz 以上共模电流无抑制。
\item 100 MHz 以上设备内部 X1/Y2 安全电容失效（寄生电感主导），共模电流穿过电源进入外部电源线。
\end{itemize}

\textbf{数量级估算}

\begin{table}[H]
\centering
\caption{电源线传导泄漏数量级估算}
\label{tab:order-of-magnitude}
\small
\begin{tabularx}{\textwidth}{lXX}
\toprule
\rowcolor{tableheader}
\textbf{参数} & \textbf{数值} & \textbf{推导依据} \\
\midrule
RPi V1 MIPI byte clock     & 51 MHz                       & 论文 Table II \\
主泄漏频点（204 / 255 MHz） & byte clock 的 4 倍 / 5 倍谐波  & 论文 Table II \\
估算 CM 电流 \(I_\mathrm{CM}\) & 1–10 \(\mu\)A（待 PhD 校准） & 类比空气场强 \(-80\) dBm + 寄生电感模型 \\
Tekbox CT 二次侧电压         & \(V_2 = Z_T \cdot I_\mathrm{CM}\)，\(Z_T \approx 5\,\Omega\) & Tekbox 数据手册 \\
等效功率 @ 50 \(\Omega\)     & \(-77 \sim -57\) dBm           & 假设 \(I_\mathrm{CM} \in [1, 10]\,\mu\)A \\
\bottomrule
\end{tabularx}
\end{table}

\textbf{关键不确定性}：\(I_\mathrm{CM}\) 的绝对值是 Phase 0 必须实测验证的核心数字。

\subsection{耦合方式权衡}
\label{sec:ch2-coupling}

\begin{table}[H]
\centering
\caption{三种耦合方式权衡（5 列）}
\label{tab:coupling-tradeoff}
\small
\begin{tabularx}{\textwidth}{lXccc}
\toprule
\rowcolor{tableheader}
\textbf{方式} & \textbf{频响 \& 隔离} & \textbf{体积} & \textbf{安全性} & \textbf{选择} \\
\midrule
宽带 CT (Tekbox TBCP2-1000) & 100 kHz – 1 GHz, \(\pm 2\) dB，磁耦合电气完全隔离 & OD \(\approx 100\) mm 夹式 & 高 & \textbf{主路} \\
HV 陶瓷电容 (1 kV/4.7 pF NP0) & 100 MHz – 数 GHz，电容短路风险，需双重串联 + LC HPF & 0805 SMD & 中 & \textbf{辅路} \\
微型 LISN (TBL5016-1) & 9 kHz – 30 MHz，自带隔离 & \(150 \times 100 \times 80\) mm & 标准 & \textbf{否决} \\
\bottomrule
\end{tabularx}
\end{table}

\textbf{选择理由}

\begin{enumerate}
\item \textbf{主路 Tekbox CT 解决 80\% 问题}：100 kHz – 1 GHz 平坦覆盖 EM Eye Table II 12/12 设备的全部主泄漏频点，NIST 可追溯校准，Phase 0 第 1 周即可投入。
\item \textbf{辅路 HV 电容补足 \(>\)500 MHz}：CT 在 500 MHz 以上磁导率下降，HV 陶瓷电容在 500 MHz – 2 GHz 段呈现低阻抗、低损耗。LDSDR 2 RX 通道允许双路同时采样、PC 端融合。
\end{enumerate}

\subsection{选定架构}
\label{sec:ch2-arch}

\paragraph{主路（CT）} Tekbox TBCP2-1000 商用宽带电流探头，100 kHz – 1 GHz 平坦 \(\pm 2\) dB，\(Z_T \approx 5\,\Omega\)（等效插损 \(-20\) dB），一次侧最大 30 A AC；L+N 同向穿芯（PE 独立走线在 CT 外），磁通 \(\Phi_\mathrm{CT} \propto I_L + I_N = 2 I_\mathrm{CM}\)，DM 抑制估算 30–40 dB；W1 制作 2 芯延长线工装，\(\delta_\mathrm{offset} \leq 0.5\,\mathrm{mm}\)。

\paragraph{辅路（HV 电容）} Murata GA355DR7GF472KY02（4.7 pF / 1 kV NP0）\(\times 2\) 串联（双重失效保护）；接入 DUT 电源线 L 对地 \(\to\) LC HPF（\(f_c = 200\) MHz）+ 50 \(\Omega\) 端接 \(\to\) SMA \(\to\) LDSDR RX2；200 MHz – 2 GHz 补足 CT 高频段；任一电容击穿，另一颗仍保持 \(> 500\) V 余量。

\paragraph{双路融合策略（LDSDR 2T2R 同时采样）}

\begin{itemize}
\item LDSDR RX1 \(\to\) CT 主路 \(\to\) LO 锁 RPi V1 @ 204 MHz（byte clock 4 次谐波）
\item LDSDR RX2 \(\to\) HV 电容辅路 \(\to\) LO 锁 @ 1.485 GHz（HDMI 扩展）或 900 MHz（EM Eye Table II 中段）
\item PC 端按 EM Eye 论文 Eq.3 做相干合成，SNR 提升估算 +3 dB \statusplan
\end{itemize}

差异化：论文用单 USRP 时分采样不同频段；本方案用 LDSDR 2 RX 同时采样。

\subsection{保护与工频隔离链}
\label{sec:ch2-protect}

完整保护链 \(\to\) 表 \ref{tab:protection-chain}（占位图 \texttt{figures/proposal\_protection\_chain.png}）。

\begin{table}[H]
\centering
\caption{保护链级联（Tekbox CT 二次侧 \(\to\) LNA 输入）}
\label{tab:protection-chain}
\small
\begin{tabularx}{\textwidth}{lXX}
\toprule
\rowcolor{tableheader}
\textbf{级} & \textbf{元件 / 配置} & \textbf{关键参数} \\
\midrule
入口 & SMA 母座 + RG-316 \(< 20\) cm + 铜罩单点接地 & 50 \(\Omega\) \\
TVS & Bourns CDSOT23-T05LC（低 C 型） & \(C < 1\) pF @ 1 GHz, V\_RWM \(\pm 5\) V, ESD \(\pm 15\) kV \\
DC 隔直 & Murata GRM18 NP0 1 nF/100 V & 50 Hz 容抗 3.18 M\(\Omega\) / 1 GHz 容抗 0.16 \(\Omega\) \\
LC HPF & 4 阶 Butterworth, \(f_c = 30\) MHz, 50 \(\Omega\) & C1=150 pF, L2=150 nH, C3=56 pF, L4=330 nH \\
输出 & ERA-4SM+ \(\times 2\) 级联板 & 50 Hz 衰减 \(> 100\) dB（目标），通带 \(< 1\) dB \\
\bottomrule
\end{tabularx}
\end{table}

\textbf{Phase 2 Route B 市电入口保护}：GDT（Bourns 2026-23-SM）\(\to\) MOV（Littelfuse V275LA20A）\(\to\) 共模扼流圈（Würth 744232222），仅插头形态需要。

\paragraph{关键工程规则} (1) 信号路径不加共模扼流圈（目标信号即共模电流本体）；(2) TVS 选低 C 型号（普通 SMAJ12A 寄生 \(\sim 200\) pF，1 GHz 容抗 0.8 \(\Omega\) 近似短路，需 \(C < 1\) pF RF 专用 TVS）；(3) 隔直电容选 NP0/C0G（\(\pm 30\) ppm/°C，Q 高，无压电；X7R 机械振动产生 \(\mu\)V 级噪声不可接受）。

\subsection{已知风险与缓解措施}
\label{sec:ch2-risks}

\begin{table}[H]
\centering
\caption{Ch2 风险清单}
\label{tab:ch2-risks}
\small
\begin{tabularx}{\textwidth}{cXXXc}
\toprule
\rowcolor{tableheader}
\textbf{\#} & \textbf{描述} & \textbf{触发条件} & \textbf{缓解措施} & \textbf{优先级} \\
\midrule
R2-1 & 实际 \(I_\mathrm{CM}\) 比假设的 1 \(\mu\)A 弱 10 dB & Phase 0 实测信号 \(< -100\) dBm 到 LNA 输入 & Phase 1 预留 PGA-103+ 第三级（+22 dB）备料 & 高 \\
R2-2 & Tekbox CT \(> 500\) MHz 段衰减大于数据手册 & VNA 实测平坦度 \(> \pm 5\) dB & 启用辅路 HV 电容 + 多频段融合软件补偿 & 中 \\
R2-3 & L+N 工装机械精度不足，DM 抑制 \(< 20\) dB & 50 Hz 残留 \(> -40\) dBm 到 LNA & 3D 打印高精度工装，\(\delta_\mathrm{offset} \leq 0.2\,\mathrm{mm}\) & 中 \\
R2-4 & 电源线驻波导致 CT 位置敏感 & 不同墙插测试 SNR 差 \(\pm 10\) dB & 统一工装距离 1.5 m，记录位置 & 低 \\
R2-5 & HV 电容辅路击穿短路 & DUT 浪涌/雷击 & 双电容串联 + GDT 一级浪涌防护 & 低 \\
R2-6 & 不同 DUT CM 阻抗差异大（20–500 \(\Omega\)） & 多设备对比部分测不到 & LDSDR AGC 76 dB 范围 + 软件自动调度 & 低 \\
\bottomrule
\end{tabularx}
\end{table}

\subsection{Phase 0 实测 Go/No-Go 决策门}
\label{sec:ch2-gonogo}

Week 1–3 内完成 5 项实验，作为是否进入 Phase 1 的硬决策门。

\paragraph{实验流程} W1 Tekbox CT VNA + 频谱仪 S21 频响 100 kHz – 1.5 GHz 扫频（要求 100 kHz – 1 GHz \(\pm 3\) dB 平坦）+ ERA-4SM+ \(\times 2\) LNA 与 LDSDR 整链注入校准（信号发生器 \(-80\) dBm \(\to\) LDSDR RSSI \(-20\) dBm，匹配 +60 dB 增益）；W2 RPi 4B + Camera V1.3 通电，Tekbox 夹电源线，频谱仪扫 50 MHz – 1 GHz 搜 30 Hz 周期性载波（目标 204 / 255 MHz）；W2–W3 LDSDR 8 MSPS IQ 采集 + Python 幅度解调 + 30 Hz 自相关 \(\to\) Tf/Tr \(\to\) 像素重排；W3 扩展到 3 COTS（Wyze/Dafang/360）。

\begin{infobox}[title=Go 条件（任一满足）]
\begin{itemize}
\item (a) RPi V1 在 204 MHz 测到 SNR \(> 15\) dB 的 30 Hz 周期信号，Python 重建结果 SSIM \(> 0.5\)
\item (b) 至少 2 个 COTS 设备在 EM Eye Table II 频点测到 SNR \(> 10\) dB
\end{itemize}
\end{infobox}

\begin{warnbox}[title=No-Go 条件（任一立即停止）]
\begin{itemize}
\item (a) 所有 5 个 DUT 在全部目标频点均测不到 SNR \(> 5\) dB（说明电源线传导通道不存在或损耗 \(> 60\) dB）
\item (b) LNA 在工作状态下饱和，无法通过 BPF 解决
\end{itemize}
\end{warnbox}

\textbf{预算前置门}：Phase 0 总预算 ¥ 8{,}860，即使 No-Go 项目总投入 \(< \) ¥ 10{,}000，无沉没成本风险。

% ============================================================================
\section{模拟前端设计}
\label{sec:ch3}
% ============================================================================

\subsection{架构：候选人 BYO 的 ERA-4SM+ \texorpdfstring{$\times 2$}{x2} 级联}
\label{sec:ch3-arch}

候选人 BYO ERA-4SM+ \(\times 2\) 级联 LNA 模块 \statusok（设计、加工与 VNA 实测均已完成，详见附录 B），模拟前端直接复用。

\begin{table}[H]
\centering
\caption{LNA 实测数据摘要（来源 \texttt{figures/lna\_era4sm\_x2\_gain.PNG}）}
\label{tab:lna-measured}
\small
\begin{tabularx}{\textwidth}{lXX}
\toprule
\rowcolor{tableheader}
\textbf{参数} & \textbf{实测值} & \textbf{备注} \\
\midrule
工作频段 & 5.6 MHz – 3 GHz & 覆盖 EM Eye Table II 12/12 设备 + HDMI 主谐波 \\
平均增益 & +28 dB @ 100 MHz & VNA 实测，网络分析仪标定后 \\
增益平坦度 & \(\pm 3\) dB @ 100 MHz – 1 GHz & 链路预算视为常数 \\
噪声系数（典型/级联） & 2.5 dB / \(\sim 3.5\) dB & 实测 NF 待 Phase 0 用噪声源标定 \\
P1dB（级联估算） & +15 dBm & 留 IIP3 余量给强干扰 \\
隔离度（失电） & \(> 50\) dB & 实测 \\
供电 & 5 V @ 60 mA（单级）/ 120 mA 总 & USB-C 可直供 \\
物理尺寸 & \(50 \times 30 \times 12\) mm & 含 SMA + 屏蔽腔 + 偏置 \\
\bottomrule
\end{tabularx}
\end{table}

\textbf{论文对标}：EM Eye 论文 Appendix H 使用 Foresight FST-RFAMP06（40 dB 单级），本方案 ERA-4SM+ \(\times 2 = 28\) dB（差 12 dB），由 LDSDR AD9363 内部 LNA + IF VGA（最高 30 dB）弥补；AD9363 级联点在 LNA 之后，对系统 NF 影响 \(< 0.1\) dB（Friis）。

\subsection{链路预算表}
\label{sec:ch3-budget}

按目标频段 200 MHz、假设 \(I_\mathrm{CM} = 1\,\mu A\) 推导，关键节点见表 \ref{tab:link-budget}。

\begin{table}[H]
\centering
\caption{链路预算（节点级）}
\label{tab:link-budget}
\small
\begin{tabularx}{\textwidth}{Xcccc}
\toprule
\rowcolor{tableheader}
\textbf{节点} & \textbf{信号 (dBm)} & \textbf{噪底 (dBm/Hz)} & \textbf{累计增益 (dB)} & \textbf{累计 NF (dB)} \\
\midrule
(1) DUT 一次侧 \(I_\mathrm{CM} = 1\,\mu\)A & — & — & — & — \\
(2) Tekbox CT 二次侧 & \(-77\) & \(-174\) & \(-20\) & 20 \\
(3) TVS + DC 隔直 + 走线 & \(-77.5\) & \(-174\) & \(-20.5\) & 20.5 \\
(4) LC HPF Stage 1 输出 & \(-78\) & \(-174\) & \(-21\) & 21 \\
(5) ERA-4SM+ 第 1 级输出 & \(-64\) & \(-171\) & \(-7\) & \textbf{3.5（主导）} \\
(6) ERA-4SM+ 第 2 级输出 & \(-50\) & \(-167\) & +7 & 3.6 \\
(7) LDSDR AD9363 入口 & \(-50\) & \(-167\) & +7 & 3.6 \\
(8) AD9363 LNA(\(-3\) dB) + IF VGA(+30 dB) & \(-23\) & \(-140\) & +34 & 3.6 \\
(9) ADC 输出动态范围 & \(-33\) dBFS & — & — & — \\
\bottomrule
\end{tabularx}
\end{table}

\paragraph{Friis NF 级联公式}

\[
NF_\mathrm{total} = NF_1 + \frac{NF_2 - 1}{G_1} + \frac{NF_3 - 1}{G_1 \cdot G_2}
\]

代入 \(NF_1 = 2.24\)，\(G_1 = 25.1\)，\(NF_2 = 2.24\)，\(G_2 = 25.1\)，\(NF_3 = 3.16\)：

\[
NF_\mathrm{total} = 2.24 + \frac{1.24}{25.1} + \frac{2.16}{630.5} = 2.29 \Rightarrow 3.6\,\mathrm{dB}
\]

\paragraph{热噪声底（8 MSPS 工作带宽）} \(P_\mathrm{noise} = -174 + 10\log(8 \times 10^6) + 3.6 = -101.4\,\mathrm{dBm}\)

\paragraph{SNR 估算} 信号 \(-77\) dBm 经 \(-20\) dB CT \(\to -97\) dBm 等效到 LNA 输入，加 28 dB LNA \(\to -69\) dBm @ LDSDR 输入，噪底折算 \(-108\) dBm，\textbf{SNR \(\approx 39\) dB}（舒适余量）。若实际 \(I_\mathrm{CM}\) 弱 30 dB，SNR 降到 9 dB 仍勉强可重建；再弱 10 dB 需启用备用 PGA-103+ 第三级（+22 dB）。

\subsection{PCB 与屏蔽}
\label{sec:ch3-pcb}

\begin{table}[H]
\centering
\caption{4 层 PCB 叠层（80 x 60 mm 模拟前端板）}
\label{tab:pcb-stackup}
\small
\begin{tabularx}{\textwidth}{cXXc}
\toprule
\rowcolor{tableheader}
\textbf{层} & \textbf{用途} & \textbf{材质} & \textbf{厚度} \\
\midrule
L1 & RF 信号顶层（LNA + 滤波） & RO4350B 局部 + FR-4 & 0.508 mm \\
L2 & 完整接地平面（绝对不切割） & FR-4 & 0.2 mm \\
L3 & 电源平面（分模拟/数字，星形互连） & FR-4 & 0.2 mm \\
L4 & 数字控制（I\textsuperscript{2}C、SPI 偏置） & FR-4 & 0.508 mm \\
\bottomrule
\end{tabularx}
\end{table}

Rogers RO4350B：\(\varepsilon_r = 3.66\,@\,1\,\mathrm{GHz}\)，损耗角正切 0.0037，50 \(\Omega\) 微带线宽 30 mil。

\paragraph{屏蔽与接地} 模拟前端板装铜罩（顶+底），罩接 L2 模拟地，单点接地到 LDSDR SOM 入口；LNA 单独子腔屏蔽避免反馈振荡；模拟/数字地通过单 0 \(\Omega\) 跳线（Star Ground）；SMA PCB 直焊，via stitching 间距 \(\leq 1/10\) 波长（@ 1 GHz 即 30 mm）。

\paragraph{电源去耦 \& EMI} LNA 偏置 5 V \(\to\) TPS7A47 LDO（4 nV/\(\sqrt{\mathrm{Hz}}\)）+ 0.1/10/100 \(\mu\)F 近端去耦；LDSDR 5 V \(\to\) TPS54320 \(\to\) 1.8 V 数字 + 独立 LDO 给 AD9363；铝合金外壳 + 铜罩屏蔽效能 40 – 60 dB @ 1 GHz。

% ============================================================================
\section{数据采集与无线传输}
\label{sec:ch4}
% ============================================================================

\subsection{SDR 选型：LDSDR 7010 rev2.1}
\label{sec:ch4-sdr}

候选人 BYO \textbf{LDSDR 7010 rev2.1} \statusok（取代原计划 ADALM-Pluto）。

\begin{table}[H]
\centering
\caption{LDSDR vs 标准 ADALM-Pluto 对比}
\label{tab:ldsdr-vs-pluto}
\small
\begin{tabularx}{\textwidth}{lXXX}
\toprule
\rowcolor{tableheader}
\textbf{项} & \textbf{LDSDR 7010（选定）} & \textbf{ADALM-Pluto} & \textbf{差异} \\
\midrule
主芯片 / RF & XC7Z010 / AD9363(70M–6G) & 同 & — \\
DDR3 & \textbf{512 MB} & 256 MB & \textbf{2 倍} IQ 缓冲 \\
网络接口 & \textbf{千兆 Ethernet + USB OTG} & 仅 USB 2.0 & \textbf{关键差异化} \\
RF 端口 & \textbf{2 TX + 2 RX（2T2R）} & 1 TX + 1 RX & \textbf{双 RX 通道} \\
扩展 I/O & 38 pin PL + 8 pin PS & 极少 & 可控外部 BPF 开关 \\
启动 & TF 卡 + 32 MB QSPI Flash & 内嵌 Flash & 调试方便 \\
候选人持有 & \statusok\ BYO，用过 OFDM + LDPC & 不持有 & \textbf{零启动成本} \\
\bottomrule
\end{tabularx}
\end{table}

\textbf{关键差异化点}：LDSDR 同款 Zynq-7010 上已实现 OFDM + LDPC 完整 HDL 收发机 BER = 0 \statusok。直接复用：(1) LDSDR Buildroot rootfs（Linux 5.15）+ iio\_attr 用户态控制 + Vivado 工程结构 + AD9363 SPI 配置流程 \statusok；(2) Phase 2 \texttt{emeye\_accel} RTL 复用 OFDM 项目 AXI-Stream / AXI-Lite / IRQ 框架；(3) XCZU3EG HLS LPR CNN（87.94\% / 675 ms）适用于 LDSDR 板上 pix2pix 加速备份方案。

\subsection{采样策略}
\label{sec:ch4-sample}

\subsubsection{频点配置}

对齐 EM Eye 论文基线 \(f_s = 8\,\mathrm{MSPS}\) IQ，中心频率锁定在 MIPI byte clock 谐波。

\begin{table}[H]
\centering
\caption{Phase 0 主目标频点}
\label{tab:phase0-freqs}
\small
\begin{tabularx}{\textwidth}{Xccc}
\toprule
\rowcolor{tableheader}
\textbf{DUT} & \textbf{byte clock (MHz)} & \textbf{目标 LO (MHz)} & \textbf{论文 SSIM} \\
\midrule
Raspberry Pi V1（论文同款） & 51 & \textbf{204 (4x)} / 255 (5x) & 0.55 / 0.61 \\
Wyze Cam Pan 2 & 222.5 & \textbf{890 (4x)} & 0.42 \\
Xiaomi Dafang & 80.5 & \textbf{322 / 890} & 0.39 \\
360 行车记录仪 M320 & 112.5 & \textbf{450} & 0.37 \\
Google Pixel 3（对照） & 128.75 & \textbf{515} & 0.34 \\
\bottomrule
\end{tabularx}
\end{table}

\subsubsection{频点跟踪算法}

\paragraph{开机扫描（Phase 1 完成后）} (1) 快扫 50 MHz – 1 GHz 步长 1 MHz，每点 100 ms 记录功率谱；(2) 粗候选 top-5 强谱线（信号 \(>\) 噪底 + 15 dB）；(3) 细识别每个候选 IQ 采集 1 s，FFT + 30 Hz 自相关识别 EM Eye 泄漏（FM 广播等无 30 Hz 周期）；(4) 锁定 + 漂移补偿（TCXO \(\pm 25\) ppm + PC 端软件 PLL \(\pm 50\) ppm）。频点漂移 \(< 50\) ppm @ 1 GHz 即 \(\pm 50\) kHz，远小于 8 MSPS 带宽。

\subsection{传输方案决策表}
\label{sec:ch4-transport}

\begin{table}[H]
\centering
\caption{原始 IQ 直传 vs 板上重建对比}
\label{tab:transport-decision}
\small
\begin{tabularx}{\textwidth}{lcccXX}
\toprule
\rowcolor{tableheader}
\textbf{方案} & \textbf{数据率} & \textbf{PC 算力} & \textbf{灵活性} & \textbf{延迟} & \textbf{选择} \\
\midrule
\textbf{v1 千兆 Eth 原始 IQ 直传} & 单 192 / 双 384 Mbps & 高（PC GPU 跑 pix2pix） & 最大 & 100 ms+ & \textbf{Phase 1} \\
\textbf{v2 板上 emeye\_accel + 8-bit 流} & 单 32 / 双 64 Mbps & 低（仅精修） & 锁定 & \(< 10\) ms & \textbf{Phase 2} \\
WiFi 6 传原始 IQ & 同上，实际 200 Mbps & 同 & 同 & + 100 ms 抖动 & 否决 \\
LDSDR USB 2.0 OTG & 限 \(\sim 10\) MSPS / 30 MB/s & — & — & — & 否决 \\
\bottomrule
\end{tabularx}
\end{table}

\textbf{选定路径}：v1 千兆 Ethernet 原始 IQ 直传（双 RX 384 Mbps，占千兆 \(\sim 48\,\%\)），Phase 2 板上 \texttt{emeye\_accel} 加速将数据率降 12 倍（384 \(\to\) 32 Mbps）。

\subsection{主控 SoC：直接用 LDSDR PS}
\label{sec:ch4-soc}

\subsubsection{决策：不外接 RPi CM4}

LDSDR 板内 Zynq-7010 PS 直接作主控，取消外接 Raspberry Pi CM4。

\begin{table}[H]
\centering
\caption{LDSDR PS 取代 CM4 的论证}
\label{tab:ldsdr-ps-vs-cm4}
\small
\begin{tabularx}{\textwidth}{lXXX}
\toprule
\rowcolor{tableheader}
\textbf{项} & \textbf{LDSDR Zynq-7010 PS（选定）} & \textbf{RPi CM4 8 GB（取消）} & \textbf{备注} \\
\midrule
CPU / 内存 & A9 双核 866 MHz / 512 MB DDR3 & A72 四核 1.5 GHz / 8 GB DDR4 & 算力够用，PC 端做重 lift \\
网络 & 千兆 Ethernet 直接 & 千兆 + WiFi 6 & WiFi 由 PC 承担 \\
OS & Buildroot Linux 5.15 & Raspbian 64-bit & 候选人已有 BSP 经验 \\
SDR 控制 & iio\_attr 直接 & USB 转串口 & 减少桥接层 \\
BOM / PCB & 0（已含） & ¥ 1000 + \(80 \times 60 \times 12\) mm & \textbf{省 ¥ 1000} \\
启动时间 & 5 s & 25 s & 单 SoC 启动快 \\
\bottomrule
\end{tabularx}
\end{table}

\textbf{唯一代价}：PS 端需跑数据搬运（从 PL 取 IQ \(\to\) 千兆 PHY DMA），Zynq-7010 PS 算力实测可承担。

\textbf{软件栈}：LDSDR PS 端（Buildroot Linux 5.15）用 \texttt{iio\_attr} 配置 AD9361 LO/增益/采样率/滤波器，libiio-dame 服务端通过千兆 socket 把 IQ 流推送到 PC，Python 端 GNURadio + pyadi-iio 实时采集；PC 端接收双 RX IQ 流（对齐时戳）\(\to\) 频点跟踪 + AGC 反馈 \(\to\) 幅度解调 + 30 Hz 自相关 \(\to\) Tf/Tr 估计 \(\to\) 像素重排 + pix2pix GAN 精修（GPU 加速）。

\subsection{板上 FPGA 加速路线图（Phase 2）}
\label{sec:ch4-accel}

LDSDR Zynq-7010 PL 实现 EM Eye 重建管线热路径加速。

\begin{table}[H]
\centering
\caption{板上加速器子模块清单}
\label{tab:accel-modules}
\small
\begin{tabularx}{\textwidth}{lXXX}
\toprule
\rowcolor{tableheader}
\textbf{模块} & \textbf{功能} & \textbf{资源估算} & \textbf{候选人能力对应} \\
\midrule
\texttt{magnitude\_cordic} & \(|I + jQ|\) CORDIC 整数幅度 & \(< 5\,\%\) LUT & OFDM BPSK/QPSK 软解调复用 \\
\texttt{autocorr\_30hz} & 30 Hz 周期自相关 + 峰值检测 & 8 \% LUT + 32 KB BRAM & OFDM 帧同步复用 \\
\texttt{decimator\_8m} & 56 \(\to\) 8 MSPS 抽取 + AAF & 12 \% LUT + 4 DSP & LDPC FIR 链路复用 \\
\texttt{iq\_to\_8bit} & 12-bit IQ \(\to\) 8-bit 解调流 & \(< 2\,\%\) LUT & — \\
\texttt{axi\_dma\_streaming} & PL \(\to\) PS Linux 内核 DMA & \(< 5\,\%\) LUT & OFDM AXI-Stream 框架复用 \\
\bottomrule
\end{tabularx}
\end{table}

\textbf{预期效果}：输出数据率 192 \(\to\) 32 Mbps（12 倍压缩）；端到端延迟 100 ms \(\to < 10\) ms；千兆占用 48 \% \(\to\) 4 \%。

\textbf{能力背书}：\texttt{emeye\_accel} RTL 首版 Vivado 工程仿真通过 \statusok（附录 B）；LDPC BER = 0 \(\to\) PS-PL 协同链路；XCZU3EG HLS LPR 675 ms \(\to\) FPGA 跑 pix2pix 级复杂度网络可行。

% ============================================================================
\section{主要器件选型 BOM}
\label{sec:ch5}
% ============================================================================

完整 BOM 见表 \ref{tab:bom-full}，涵盖 Route A Phase 1 单台原型所需全部器件。\textbf{BYO 标注的元件不计入 BOM 成本}。

\begin{table}[H]
\centering
\caption{Route A Phase 1 完整 BOM}
\label{tab:bom-full}
\footnotesize
\begin{tabularx}{\textwidth}{clXrX}
\toprule
\rowcolor{tableheader}
\textbf{\#} & \textbf{模块} & \textbf{主选型号 / 替代} & \textbf{单价 (¥)} & \textbf{选择理由} \\
\midrule
1 & 耦合主路 & Tekbox TBCP2-1000 / Pearson 411 & 4{,}500 & NIST 校准 + Phase 0 第 1 周可用 \\
2 & 耦合辅路 \(>500\) MHz & Murata GA355 4.7 pF/1 kV NP0 \(\times 2\) 串联 & 30 & 双重失效保护 + RF 低损耗 \\
3 & 保护 TVS & Bourns CDSOT23-T05LC / TI TPD3E001 & 8 & \(C < 1\) pF @ 1 GHz，RF 透明 \\
4 & 保护 DC 隔直 & Murata GRM18 NP0 1 nF/100 V \(\times 2\) & 6 & NP0 \(\pm 30\) ppm，无压电 \\
5 & 保护 GDT（Phase 2） & Bourns 2026-23-SM & 12 & Phase 1 不需要（USB 供电） \\
6 & 保护 MOV（Phase 2） & Littelfuse V275LA20A & 8 & Phase 1 不需要 \\
7 & 保护 共模扼流（Phase 2） & Würth 744232222 & 25 & Phase 1 不需要 \\
8 & 滤波 Stage 1 必备 & Murata 150/56 pF NP0 + Coilcraft 0603HP & 16 & LC HPF 4 阶 Bw 50\(\Omega\) fc=30M \\
9 & LNA 主放大 & \textbf{ERA-4SM+ \(\times 2\) 级联（BYO）} & \textbf{0} & VNA 实测 +28 dB，5.6M–3G \\
10 & LNA 备用第三级 & Mini-Circuits PGA-103+ & 100 & Phase 0 测出 SNR 不够时启用 \\
11 & 滤波 Stage 2（可选） & Murata SAFEB/SAFEA/SAFFB SAW & 250 & Phase 0 实测决定 \\
12 & RF 开关（可选） & Peregrine PE42423 SP4T & 100 & XCZU3EG 项目经验 \\
13 & SDR & \textbf{LDSDR 7010 rev2.1（BYO）} & \textbf{0} & BYO + 千兆 + 2 RX \\
14 & 主控 & \textbf{LDSDR Zynq-7010 PS（BYO）} & \textbf{0} & Buildroot + iio\_attr \\
15 & 电源 模拟 LDO & TI TPS7A47 / LT3045 & 20 & 4 nV/\(\sqrt{\mathrm{Hz}}\) LNA 净化电源 \\
16 & 电源 数字开关 & TI TPS54320 / LM5085 & 15 & USB 5 V \(\to\) 1.8 V 数字 \\
17 & 电源 USB-C 输入 & Maxim MAX14778 + ESD7104 & 30 & USB-C 5 V/3 A + ESD 保护 \\
18 & PCB 模拟前端 & 4 层 \(80 \times 60\) mm，RO4350B 局部 & 400 & 小批量打样 5 块，含贴片 \\
19 & 接插件 + 屏蔽 & SMA \(\times 3\) + RJ45 + USB-C + 铜罩 & 60 & 含 EMI 衬垫 \\
20 & 网线 & CAT6 1 m & 10 & LDSDR \(\to\) PC \\
21 & 外壳 Phase 1 & 3D 打印 ABS，\(100 \times 60 \times 30\) mm & 150 & 快速迭代 2 次 \\
22 & 外壳 Phase 2 & CNC 铝合金，带 EMI 衬垫 & 800 & Phase 2 演进至插头形态 \\
\bottomrule
\end{tabularx}
\end{table}

\paragraph{Phase 1 BOM 净成本汇总}

\begin{table}[H]
\centering
\caption{BOM 汇总}
\label{tab:bom-summary}
\small
\begin{tabularx}{\textwidth}{lrX}
\toprule
\rowcolor{tableheader}
\textbf{类别} & \textbf{金额 (¥)} & \textbf{备注} \\
\midrule
主路 + 辅路耦合 & 4{,}530 & Tekbox 大头 + HV 电容 \\
保护链 Phase 1 必备 & 14 & TVS + DC 隔直 \\
滤波 Stage 1 必备 & 16 & LC HPF \\
LNA + SDR + 主控（BYO） & 0 & 候选人 BYO \\
备用 LNA 第三级 + Stage 2 BPF & 450 & Phase 0 触发才启用 \\
电源 + USB-C & 65 & 含 ESD \\
PCB / 接插件 / 网线 / 外壳 & 620 & 4 层 + RO4350B + 3D 打印 \\
\textbf{基线 BOM 合计} & \textbf{5{,}245} & 含 Tekbox \\
\textbf{不含 Tekbox（实验室借用）} & \textbf{745} & 仅工程化部分 \\
\bottomrule
\end{tabularx}
\end{table}

% ============================================================================
\section{整体成本估算与交付计划}
\label{sec:ch6}
% ============================================================================

\subsection{BOM 成本}
\label{sec:ch6-bom}

Phase 1 单台原型 BOM 合计 \textbf{¥ 5{,}245}（含 Tekbox CT）或 \textbf{¥ 745}（实验室借用 Tekbox）。BYO 资产（LDSDR + ERA-4SM+ \(\times 2\) LNA + Vivado/HLS/Buildroot 工具链）估值 ¥ 15{,}000–20{,}000，\textbf{不计入项目预算}。

\subsection{NRE（非经常性工程）成本}
\label{sec:ch6-nre}

\begin{table}[H]
\centering
\caption{Phase 0 + 1 + 2 NRE 成本汇总}
\label{tab:nre-summary}
\small
\begin{tabularx}{\textwidth}{lXrX}
\toprule
\rowcolor{tableheader}
\textbf{阶段} & \textbf{项} & \textbf{金额 (¥)} & \textbf{说明} \\
\midrule
\textbf{Phase 0} & Tekbox TBCP2-1000（若需采购） & 4{,}500 & 实验室有则为 0 \\
                & RPi 4B \(\times 2\) + Camera V1.3 \(\times 2\) & 1{,}160 & 论文同款 DUT \\
                & COTS 摄像头 \(\times 4\)（Wyze/Dafang/360/Pixel） & 1{,}200 & 多设备验证 \\
                & SMA 同轴 + 转接 + 杂项 & 800 & 调试耗材 \\
                & 隔离变压器 1:1 220V/1kVA & 1{,}200 & 实验室有则为 0 \\
                & Phase 0 最坏 / 实际（借用） & \textbf{8{,}860 / 3{,}160} & 仅 DUT + 耗材 \\
\textbf{Phase 1} & PCB 打样（2 次 \(\times 5\) 块） & 1{,}500 & RO4350B 顶层 \\
                & 3D 外壳打样（2 次） & 300 & ABS \\
                & 测试设备时间（VNA / 频谱仪） & 3{,}000 & 实验室借用，机时 \\
                & 调试耗材 & 300 & \\
                & Phase 1 小计 & \textbf{5{,}100} & \\
\textbf{Phase 2} & 6 层集成 PCB（1 次） & 2{,}500 & + IEC 间距 \\
                & CNC 铝合金外壳 & 800 & 替代 3D 打印 \\
                & 自研小型化 CT 探索 & 600 & Fair-Rite 磁芯 + 工具 \\
                & 备用 LDSDR & 1{,}500 & 防硬件故障 \\
                & 频谱仪扩展头（\(> 3\) GHz） & 600 & 备用 \\
                & Phase 2 小计 & \textbf{6{,}000} & \\
\textbf{NRE 总计} & & \(\sim\) \textbf{14{,}260} & 实际可压到 \(\sim 5{,}300\) \\
\bottomrule
\end{tabularx}
\end{table}

\subsection{工时与里程碑}
\label{sec:ch6-milestones}

\textbf{1 RA \(\times 6\) 个月全职}（132 个工作日）。

\paragraph{Phase 0（W1–W3）} Tekbox CT 标定 + RPi V1 复现 + COTS 设备扫描 + Go/No-Go 决策门；交付 Phase 0 实验报告 + Phase 1 设计修正建议。

\paragraph{Phase 1（W4–W15）} W4–W7 模拟前端板设计 + PCB 打样 v1 + LNA/滤波链联调；W8–W11 LDSDR Buildroot rootfs + emeye 算法 Python + 多频段双 RX 融合；W12–W14 PCB 打样 v2 + 3D 外壳 + 整机集成 + 多 DUT 验证；W15 Phase 1 评审，出 v1 完整工程原型。

\paragraph{Phase 2（W16–W24）} W16–W19 \texttt{emeye\_accel} RTL 板上集成 + PS-PL 联调（OFDM 经验直接迁移）；W20–W22 6 层集成 PCB + CNC 外壳 + 自研小型化 CT 验证；W23–W24 多设备 + 多场景验证 + 终版报告。


\subsection{Gantt 表（文字版）}
\label{sec:ch6-gantt}

\begin{lstlisting}[basicstyle=\ttfamily\scriptsize]
W:   1  2  3 | 4  5  6  7  8  9 10 11 12 13 14 15 | 16 17 18 19 20 21 22 23 24
Phase 0  ========                                                              (可行性 + Go/No-Go)
Phase 1            ===========                                                  (模拟前端 PCB v1 + LNA 联调)
                              ===========                                       (LDSDR rootfs + 算法 Python)
                                          ===========                           (PCB v2 + 外壳 + 整机集成)
                                                       =                        (Phase 1 评审)
Phase 2                                                  ================       (板上 emeye_accel + 自研小 CT)
                                                                         =====  (多设备验证 + 终版报告)
\end{lstlisting}

\subsection{总预算汇总与应急储备}
\label{sec:ch6-budget}

\begin{table}[H]
\centering
\caption{6 个月项目总预算（含应急储备）}
\label{tab:total-budget}
\small
\begin{tabularx}{\textwidth}{lrX}
\toprule
\rowcolor{tableheader}
\textbf{类别} & \textbf{金额 (¥)} & \textbf{备注} \\
\midrule
BYO 资产（LDSDR + LNA + 工具链） & 0（估值 15–20 k） & 不计预算 \\
Phase 0 采购（实际，借用大头） & 3{,}160 & 完整 \(\sim 8{,}860\) \\
Phase 1 BOM（单台原型） & 5{,}245 & 含 Tekbox \\
Phase 1 NRE & 5{,}100 & \\
Phase 2 采购 + NRE & 6{,}000 & 含备用 LDSDR + 自研 CT \\
应急预备金（PGA-103+ + 辅路 + 备板） & \(\sim 6{,}350\) & 占总硬件 \(\sim 30\,\%\) \\
项目硬件 + NRE 合计 & \(\sim\) \textbf{19{,}505} & 最低 \(\sim 14{,}000\) \\
RA 津贴（¥ 2{,}000 / 月 \(\times 6\)） & 12{,}000 & 校内 RA 标准 \\
\textbf{项目总成本} & \(\sim\) \textbf{31{,}500} & 区间 26{,}000 – 36{,}000 \\
\bottomrule
\end{tabularx}
\end{table}

\textbf{预算合理性}：BYO 资产省去 SDR + LNA 大头；Phase 0 Go/No-Go 限制沉没成本（即使 No-Go 投入 \(<\) ¥ 10{,}000 即可终止）；LISN/频谱仪/VNA 实验室公共资源借用。

% ============================================================================
\appendix
% ============================================================================

\section{候选人工程经验与能力支撑}
\label{sec:appA}

候选人项目清单（7 项，均含实测数字与代码归档）。

\subsection{OFDM + LDPC PlutoSDR Zynq-7010 全链路收发机}
\label{sec:appA-1}

平台 PlutoSDR / LDSDR（\texttt{xc7z010clg225-1}）/ 单人独立 PHY+Baseband HDL+板级联调 / 8 个月。

\paragraph{技术 \& 结果} 完整 OFDM 链路 HDL（FFT/IFFT、循环前缀、导频插入、信道估计、均衡）+ LDPC 编解码器 HDL（min-sum 迭代 + 比特节点并行）+ 时频同步（Schmidl-Cox + 导频辅助 CFO）+ Zynq-7010 PL 综合时序闭环 30.72 MHz + GNURadio baseband 实时收发。\textbf{板级实测 BER = 0} \statusok（SNR \(\geq 10\) dB, \(10^6\) bit 零误码）；Vivado LUT 38\% / BRAM 52\% / DSP 24\%；吞吐 1.5 Mbps（QPSK, 1/2 LDPC）。

\paragraph{对本提案支撑} Pluto/LDSDR Zynq-7010 内部 HDL 修改 + bitstream 板级验证 \(\to\) "板上 FPGA 加速"路径（详见 §A.8）。工具链 Vivado HLS 2020.2 / Verilog / GNURadio 3.8。

\subsection{XCZU3EG PL CNN 车牌识别端到端部署}
\label{sec:appA-2}

平台 XCZU3EG MPSoC（FZ3A 板，A53\(\times\)4 + R5 + PL）/ 单人 CNN+HLS+PL+Petalinux / 6 个月。

\paragraph{技术 \& 结果} LPRNet（CTC loss，0.5 M 参数）+ Vivado HLS（line-buffer 行流水，激活 LUT 化，INT8）+ PL/PS 协同（AXI4-Stream DMA + 中断）+ Conv-BN-ReLU pipeline II=1 时序闭环 200 MHz + Petalinux。\textbf{字符识别 87.94\%} \statusok（2000 张测试集），\textbf{端到端 675 ms}，PL LUT 71\% / DSP 83\% / BRAM 64\%。

\paragraph{对本提案支撑} PS/PL 协同 \(\to\) \texttt{emeye\_accel} PL IP + PS Linux 通信。仓库 \href{https://github.com/stongry/FPGA-ZYNQ}{github.com/stongry/FPGA-ZYNQ}。

\subsection{QSM368ZP RK3568 Ubuntu 22.04 移植与 RKNN NPU 部署}
\label{sec:appA-3}

平台 RK3568（A55\(\times\)4 + Mali-G52 + 0.8 TOPS NPU）/ QSM368ZP 工业模块 / 单人 Buildroot \(\to\) Ubuntu 22.04 全量移植 + RKNN / 4 个月。

\paragraph{技术 \& 结果} U-Boot + Linux 5.10（Rockchip BSP）+ DTS 适配 + Ubuntu 22.04 server rootfs + systemd + DSI 触摸驱动 + RKNN Toolkit 2.3.2 + RetinaFace USB camera \(\to\) RGA \(\to\) NPU \(\to\) NMS 全链路。\textbf{生产级嵌入式 Linux 已交付客户} \statusok；RetinaFace 端到端 41.7 fps @ RK3568 NPU，单帧 15 ms。

\paragraph{对本提案支撑} Phase 2 板上 PS Linux 侧 TCP forwarder + IQ 落盘服务（设备树 + udev + systemd）+ NPU RKNN 流水线（如需 SSIM offload）。

\subsection{EC800M LTE Cat-1 模块语音 AI 对话集成}
\label{sec:appA-4}

平台 Quectel EC800M（MDM9205，LTE Cat-1 + TTS/Audio）/ 单人 QuecPython + LTE 拨号 + 云端语音 AI / 3 个月。

\paragraph{技术 \& 结果} QuecPython 纯 Python 业务栈 + LTE Cat-1 PPP 拨号 + TCP/TLS 长连接到云端 ASR/LLM/TTS + 板载 audio codec 16 kHz PCM + 流式 TTS 半双工抢答 + AT 命令 OTA。\textbf{出货级语音 AI 模块批量出货} \statusok；SIM 注网到首字延迟 \(\approx 800\) ms。

\paragraph{对本提案支撑} 出货级嵌入式产品交付 + LTE 集成 + 云端长连接 + OTA。仓库 \href{https://github.com/stongry/2026yiyuan}{github.com/stongry/2026yiyuan}。

\subsection{Smart Home Slint UI 全量移植（Flutter \texorpdfstring{$\to$}{->} Slint）}
\label{sec:appA-5}

平台 RK3588 显控面板（LubanCat-5，A76\(\times\)4 + A55\(\times\)4，Mali-G610）/ 单人 Flutter \(\to\) Slint 全量重写 / 2 个月。

\paragraph{技术 \& 结果} Slint UI（Rust）+ 自定义 widget + 交叉编译 \texttt{aarch64-unknown-linux-gnu} + Wayland + EGL 后端 + tokio 异步 + Slint 属性桥接 + GPU 状态切换优化。\textbf{62.97 fps 达 60 Hz vsync 硬件上限} \statusok；二进制体积比 Flutter 下降 \(\sim 70\)\%。

\paragraph{对本提案支撑} 上位机可视化（帧显示 + SSIM 曲线 + 频谱图）Rust + Slint 高性能 UI。

\subsection{RetinaFace NPU 性能 Benchmark（RK3568 vs RK3588）}
\label{sec:appA-6}

对比 RK3568（QSM368ZP）vs RK3588（LubanCat-5）/ 单人 benchmark 方法学 + A/B 实测 / 2 周。

\paragraph{技术 \& 结果} 严格 A/B 控制（同 ONNX + 同 RKNN + 同测试集 + 同推理代码）+ 测量分离（NPU kernel / PS\(\leftrightarrow\)NPU DMA / CPU 前后处理）+ 每平台 1000 帧 + p50/p95/p99 + CPU governor 固定 + 单线程绑定。\textbf{NPU 单帧} RK3568 = 15 ms，RK3588 = 7.75 ms（\(\sim\) 2\(\times\)）；\textbf{端到端 fps} 20.7 / 41.7 \statusok。

\paragraph{对本提案支撑} Phase 0/1 A/B 实测严谨方法学 \(\to\) 完整实验设计 + 数据采集 + 统计报告。

\subsection{本提案 Pre-Phase 0 已完成工作}
\label{sec:appA-7}

平台 LDSDR + 自研耦合电路 + PC 仿真 / 6 周 / 自费预研。关键交付件：

\begin{itemize}
\item 5 个 Python 端到端仿真：EM Eye 信号建模、HDMI TMDS bit-level、TCP demo，帧级 SSIM = 0.99 \statusok
\item 4 个 RTL testbench 全 PASS：Verilator + Vivado Simulator 双验证 \statusok
\item LDSDR 真目标 bitstream 生成 + 烧录 + RF 实测：Zynq-7010 启动正常，采真实 RF 数据 \statusok
\end{itemize}

详见附录 B。

\subsection{能力 — 提案任务映射表}
\label{sec:appA-8}

\begin{table}[H]
\centering
\caption{能力 — 提案任务映射}
\label{tab:capability-mapping}
\footnotesize
\begin{tabularx}{\textwidth}{XXXc}
\toprule
\rowcolor{tableheader}
\textbf{提案任务（按时间序）} & \textbf{所需核心能力} & \textbf{对应项目} & \textbf{经验等级} \\
\midrule
Phase 0 W1：基线 EM Eye demo 复现 + Python 仿真 & Python + 信号处理 + 端到端仿真 & §A.7 & 熟练（已完成） \\
Phase 0 W2：LISN + LDSDR 烧 OFDM bitstream 实测电源线 SNR & Pluto/LDSDR HDL + bitstream + RF 采集 & §A.1 + §A.7 & 熟练 \\
Phase 0 W3：耦合电路 PCB 打样 + 简单 EMI 测试 & PCB layout + RF 走线（基础） & §A.1 SDR 板级 & 入门—有经验 \\
Phase 0 W4：基线性能 A/B benchmark & 严格 A/B 方法学 & §A.6 & 熟练 \\
Phase 1 W5–W6：emeye\_accel IP 顶层设计 & Vivado HLS + RTL 多级流水线 & §A.2 & 熟练 \\
Phase 1 W7–W8：emeye\_accel 集成进 LDSDR bitstream & Zynq-7010 PL 综合时序闭环 & §A.1 & 熟练 \\
Phase 1 W8：emeye\_accel A/B 性能对比 & A/B benchmark + 统计 & §A.6 & 熟练 \\
Phase 2 W9–W10：板上 PS Linux TCP forwarder / 落盘 & 嵌入式 Linux + systemd + 驱动 & §A.3 & 熟练 \\
Phase 2 W11：（可选）SSIM offload 到板上 NPU & NPU 生态 + 模型转换 & §A.3 & 熟练 \\
Phase 2 W12：上位机可视化重建 & Rust + Slint 或同级 UI & §A.5 & 熟练 \\
Phase 3（可选）：远程 telemetry & LTE 集成 + 云端长连接 & §A.4 & 有经验 \\
全周期：版本管理 / CI / 实验记录 & Git + Python + bash & §A.1—§A.7 & 熟练 \\
\bottomrule
\end{tabularx}
\end{table}

\footnotesize\textit{经验等级：熟练 = 独立完成过端到端项目并交付板级/客户出货结果；有经验 = 参与过相关项目部分模块；入门 = 理解原理可指导下完成。}\normalsize

\subsection{工具链与开发环境}
\label{sec:appA-9}

\begin{table}[H]
\centering
\caption{工具链与开发环境}
\label{tab:toolchain}
\small
\begin{tabularx}{\textwidth}{lXc}
\toprule
\rowcolor{tableheader}
\textbf{类别} & \textbf{工具/平台} & \textbf{熟练度} \\
\midrule
FPGA 综合 & Vivado / Vivado HLS 2020.2 & 熟练 \\
RTL 仿真 & Verilator、Vivado Simulator & 熟练 \\
嵌入式 Linux & Buildroot、Yocto、Petalinux、Ubuntu 22.04 BSP & 熟练 \\
NPU 推理 & RKNN Toolkit 2.3.2、Vitis AI 1.4 & 熟练 \\
编程语言 & Python、C/C++、Rust、Verilog/SystemVerilog、QuecPython & 熟练 \\
UI 框架 & Slint (Rust)、Flutter & 熟练 / 有经验 \\
通信协议 & LTE Cat-1 / PPP / TCP/TLS / AXI4-Stream / DMA & 熟练 \\
实验仪器 & LISN、频谱仪、示波器、SDR（PlutoSDR/LDSDR） & 熟练 \\
版本管理 & Git、GitHub、Linux shell 脚本 & 熟练 \\
\bottomrule
\end{tabularx}
\end{table}

合计 6 个已交付项目 + 1 项 Pre-Phase 0，约 23 个月有效工时，覆盖 Phase 0/1/2 关键技术栈。

\section{Pre-Phase 0 已完成工作（2026-05-12 至 2026-05-14）}
\label{sec:appB}

公开 GitHub 仓库与本地工程目录可复现、可审计。

\subsection{一句话总览}
\label{sec:appB-0}

6 天工程冲刺完成 5 类产物：(1) 2 份 Python 仿真（EM Eye SSIM 0.98 + HDMI TEMPEST SSIM 0.9907）+ 1 份 LDSDR \(\to\) 主机 TCP demo；(2) 4 个可综合 Verilog 模块 + 4 个 testbench 全 PASS（Icarus Verilog，无 Vivado license 依赖）；(3) 完整 Vivado 流程到 \texttt{write\_bitstream}，目标真器件 XC7Z010-CLG400-2（LDSDR rev2.1），\texttt{.bit} MD5 已固定；(4) 真板烧录 + AD9361 RF 实测（千兆 GbE 在线，FPGA Manager \texttt{operating}，RX\_LO 204 MHz, RSSI 93.75 dB, 16 KB IQ, dmesg 零错误）；(5) Phase 2 规划文档 1{,}096 行（AXI-Stream tap + PS TCP forwarder 骨架）。工程产物合计约 6{,}000 行。

\subsection{端到端 Python 仿真：EM Eye 论文复现到 HDMI TEMPEST 泛化}
\label{sec:appB-1}

Python 端到端链路打通后，FPGA RTL 输出有 golden reference 对比。

\subsubsection{EM Eye 论文方法完整复现（Long et al., NDSS 2024）}

文件 \texttt{simulation/emeye\_simulation.py}（501 行）：

\begin{table}[H]
\centering
\caption{EM Eye 仿真组件}
\label{tab:emeye-sim}
\small
\begin{tabularx}{\textwidth}{lXX}
\toprule
\rowcolor{tableheader}
\textbf{组件} & \textbf{实现内容} & \textbf{关键参数} \\
\midrule
信号源 & RPi V1（OV5647）BT.656-like 像素流 & 30 fps, 640\(\times\)480, MIPI byte clock 谐波 \\
信道模型 & 像素 \(\to\) byte clock 谐波 \(\to\) IQ 基带 & 加性高斯噪声 + 多径 \\
解调链 & 复数 IQ \(\to\) 幅度 \(\to\) 抽取 & 抽取因子 8 \\
帧同步 & 自相关搜索 + 滞后阈值 & 与论文 §4.2 对齐 \\
重建质量 & 与源图像 SSIM & \(\sim 0.98\) \\
\bottomrule
\end{tabularx}
\end{table}

输出 \texttt{simulation/output/}：\texttt{source\_image.png}、\texttt{reconstructed.png}、\texttt{simulation\_result.png}（三幅一组）。

\subsubsection{HDMI TEMPEST 扩展（原创泛化验证）}

文件 \texttt{simulation/hdmi\_simulation.py}（423 行）证明：论文方法是通用 EM 侧信道重建平台。

\begin{table}[H]
\centering
\caption{EM Eye 与 HDMI TEMPEST 场景对比}
\label{tab:hdmi-vs-emeye}
\small
\begin{tabularx}{\textwidth}{lXX}
\toprule
\rowcolor{tableheader}
\textbf{参数} & \textbf{EM Eye 场景} & \textbf{HDMI TEMPEST 场景} \\
\midrule
帧率 / 像素时钟 & 30 fps / 12 MHz MIPI byte clk & \textbf{60 fps / 148.5 MHz}（1080p60） \\
同步信号 / 分辨率 & 隐式（行长度推断）/ 640\(\times\)480 & \textbf{显式 H/V\_SYNC} / 1080p \(\to\) 160\(\times\)90 \\
重建 SSIM / Correlation & \(\sim 0.98\) / — & \textbf{0.9907 / 0.9976} \\
端到端运行时间 & — & \textbf{1.2 s} \\
\bottomrule
\end{tabularx}
\end{table}

输出 \texttt{simulation/output/hdmi\_simulation\_result.png}。直接支撑本提案 Phase 4"多模态侧信道"路线图。

\subsubsection{LDSDR \(\to\) 主机 TCP 网络管线 demo}

文件 \texttt{simulation/network\_demo.py}（405 行）：Python 模拟 FPGA 加速器输出 + TCP 传输 + 主机重建。Phase 2 真硬件 bring-up 时把 Python mock 换成真 socket 即可。与 §B.5 PS 用户态 TCP forwarder 骨架协议兼容（同帧头 + 同 byte order）。

\subsection{FPGA 加速器 RTL 与单元测试}
\label{sec:appB-2}

\subsubsection{RTL 模块清单}

\begin{table}[H]
\centering
\caption{RTL 模块清单}
\label{tab:rtl-modules}
\small
\begin{tabularx}{\textwidth}{Xccc}
\toprule
\rowcolor{tableheader}
\textbf{文件} & \textbf{行数} & \textbf{功能} & \textbf{DSP/BRAM} \\
\midrule
\texttt{fpga\_accel/rtl/magnitude\_jpl.v} & 117 & JPL 近似幅度 \(|I+jQ| \approx \alpha \cdot \max + \beta \cdot \min\) & 0 / 0 \\
\texttt{fpga\_accel/rtl/cic\_decimator.v} & 86 & 1 阶 CIC boxcar 抽取（8x） & 0 / 0 \\
\texttt{fpga\_accel/rtl/frame\_sync.v} & 142 & 帧同步 FSM（平均 + 滞后阈值） & 0 / 1 \\
\texttt{fpga\_accel/rtl/emeye\_accel\_top.v} & 282 & 顶层 + 双通道 + AXI-Stream FIFO & 0 / 0 \\
\textbf{RTL 小计} & \textbf{627} & & \textbf{0 DSP} \\
\bottomrule
\end{tabularx}
\end{table}

\texttt{magnitude\_jpl.v} 核心是 \textbf{JPL 系数 0.375 = 1/4 + 1/8 推导}，幅度计算可用 2 次右移 + 1 次加法完成，\textbf{完全无乘法器}：

\begin{lstlisting}[language=verilog]
// magnitude_jpl.v 核心片段(伪)
wire [W-1:0] max_ab = (abs_i > abs_q) ? abs_i : abs_q;
wire [W-1:0] min_ab = (abs_i > abs_q) ? abs_q : abs_i;
// alpha ~= 1.0, beta = 0.375 = 0x60 in Q8 -- 用移位实现
wire [W+1:0] mag = max_ab + (min_ab >> 2) + (min_ab >> 3);
\end{lstlisting}

\subsubsection{单元测试结果（Icarus Verilog 13.0）}

\begin{table}[H]
\centering
\caption{单元测试结果}
\label{tab:unit-tests}
\small
\begin{tabularx}{\textwidth}{lXcX}
\toprule
\rowcolor{tableheader}
\textbf{Testbench} & \textbf{文件 / 用例} & \textbf{结果} & \textbf{关键指标} \\
\midrule
\texttt{tb\_magnitude\_jpl} & tb\_magnitude\_jpl.v(177) / 60 & \textbf{60/60 PASS} & 平均误差 4.29\%，峰值 6.80\% \\
\texttt{tb\_cic\_decimator} & tb\_cic\_decimator.v(222) / 11 & \textbf{11/11 PASS} & 抽取因子准确，无样本溢出 \\
\texttt{tb\_frame\_sync} & tb\_frame\_sync.v(176) / 5 & \textbf{5/5 PASS} & 正确触发 \texttt{frame\_start}，正确忽略短 blanking \\
\texttt{tb\_emeye\_accel} & tb\_emeye\_accel.v(222) / 顶层 & \textbf{PASS} & 1156 AXI-Stream 输出，4 个 \texttt{frame\_starts}，零样本丢失 \\
\bottomrule
\end{tabularx}
\end{table}

\texttt{tb\_emeye\_accel} 在仿真期发现并修复了一个真实 RTL bug：顶层 round-robin 输出无缓冲导致样本丢失。修复方案是加 pending buffer + 优先级回退。

\subsubsection{设计推导文档}

\texttt{fpga\_accel/doc/fpga\_derivations.md}（529 行）系统记录 7 类公式推导：JPL 系数误差界；1 阶 CIC sinc 频响零点（1/2/3/4 MHz @ Fs=8M）；帧同步 FSM 滞后阈值；流水线时延预算（最优 7 cycle / 最坏 14 cycle）；资源预估（\(\sim 600\) LUT / 476 FF / 1 BRAM / 0 DSP）；跨时钟域 metastability 风险；7 个学长可能追问的 Q\&A 自查清单（\texttt{Why JPL?}、\texttt{Why 1 阶 CIC?}、\texttt{Why no DSP?} 等）。

\subsection{Vivado 工程化与真目标 bitstream 生成}
\label{sec:appB-3}

\textbf{目标器件：\texttt{xc7z010clg400-2}（LDSDR rev2.1 真硬件，非 PlutoSDR z020）}。

\subsubsection{完整流程时间线}

\begin{table}[H]
\centering
\caption{Vivado 流程时间线}
\label{tab:vivado-timeline}
\small
\begin{tabularx}{\textwidth}{lcX}
\toprule
\rowcolor{tableheader}
\textbf{阶段} & \textbf{状态} & \textbf{时间戳} \\
\midrule
synth\_2（z010 license 获取） & \statusok & 2026-05-13 23:17 \\
opt\_design & \statusok & DRC \textbf{0 errors} \\
place\_design / route\_design & \statusok & — \\
write\_bitstream & \statusok & 2026-05-13 23:24，\texttt{Bitgen Completed Successfully} \\
bitstream 产物 & \texttt{bitstream/ldsdr\_2tr\_emeye.bit} & 2.0 MB \\
MD5（固定） & \texttt{01a664a91e77f5477190d4f50a6de2a6} & \\
\bottomrule
\end{tabularx}
\end{table}

全流程 \textbf{7 分钟}，本地 Vivado 无 license 冲突跑完。

\subsubsection{XC7Z010 资源利用率实测}

\begin{table}[H]
\centering
\caption{XC7Z010 资源利用率}
\label{tab:resource-util}
\small
\begin{tabularx}{\textwidth}{lccX}
\toprule
\rowcolor{tableheader}
\textbf{资源} & \textbf{已用 / 总量} & \textbf{占比} & \textbf{备注} \\
\midrule
Slice LUT & 2{,}420 / 17{,}600 & \textbf{13.75\%} & 86\%+ 余量 \\
LUT as Logic & 2{,}116 / 17{,}600 & 12.02\% & \\
LUT as Memory & 304 / 6{,}000 & 5.07\% & \\
Slice Register (FF) & 4{,}074 / 35{,}200 & \textbf{11.57\%} & \\
F7 / F8 Mux & 84 / 32 & \(< 1\,\%\) & \\
\textbf{DSP} & \textbf{0} / 80 & \textbf{0\%} & JPL 完全无乘法器 \\
BRAM & 1 / 60 & 1.67\% & 平均窗口环形缓冲 \\
\bottomrule
\end{tabularx}
\end{table}

\textbf{关键结论}：86\%+ 资源余量可容纳 3-4 路并行处理通道，多频段融合可行。

\subsubsection{时序结果}

\begin{table}[H]
\centering
\caption{时序结果}
\label{tab:timing}
\small
\begin{tabularx}{\textwidth}{XccX}
\toprule
\rowcolor{tableheader}
\textbf{时钟域} & \textbf{WNS} & \textbf{状态} & \textbf{备注} \\
\midrule
\texttt{clk\_fpga\_0}（8 MHz, emeye\_accel 域） & \textbf{+3.05 ns} & \statusok\ 满足 & 候选人新加电路时序全部 clean \\
\texttt{rx\_clk \(\to\) clk\_fpga\_0}（跨时钟域） & \textbf{\(-2.985\) ns} & \statuspartial\ 123 failing & 原 ad9361 模板就有的问题 \\
\bottomrule
\end{tabularx}
\end{table}

跨时钟域 WNS 负值源于 LDSDR\_2TR 模板将 \texttt{rx\_clk} 直接打到 AXI 寄存器读侧（候选人 RTL 未引入）。Phase 1 W1 缓解措施二选一：

\begin{itemize}
\item \textbf{方案 A}：\texttt{set\_false\_path -from [get\_clocks rx\_clk] -to [get\_clocks clk\_fpga\_0]} + 加 2-FF 同步器
\item \textbf{方案 B}：在跨域路径上插 async FIFO（Xilinx FIFO Generator IP）
\end{itemize}

选定方案 A（AXI 寄存器读侧仅要求最终一致性）。

\subsubsection{Bitstream \(\to\) BIN 后处理工具}

文件 \texttt{bitstream/bit\_to\_bin.py}（53 行，纯 Python 无外部依赖）。实现：找到 Xilinx \texttt{.bit} 文件里的 sync word \texttt{0xAA995566}；后面所有 word \textbf{32-bit byte swap}（Linux \texttt{fpga\_manager} 接受 big-endian 字流）；输出 raw \texttt{.bin} 给 \texttt{/sys/class/fpga\_manager/fpga0/firmware}。产物 \texttt{bitstream/ldsdr\_2tr\_emeye\_safe.bin}，MD5 \texttt{b294f2a3ed77adffc3bebaadb6c4e538}。

\subsection{真 LDSDR 板烧录与 AD9361 RF 链路实测}
\label{sec:appB-4}

\subsubsection{板级环境}

\begin{table}[H]
\centering
\caption{板级环境}
\label{tab:board-env}
\small
\begin{tabularx}{\textwidth}{lX}
\toprule
\rowcolor{tableheader}
\textbf{项目} & \textbf{实测值} \\
\midrule
板卡 & LDSDR rev2.1（XC7Z010 + AD9361） \\
网络 & 千兆 GbE \(\to\) 本机 LAN，IP \textbf{192.168.3.10} \\
登录 & SSH \texttt{root/analog}（沿用 PlutoSDR 默认凭据） \\
板上系统 & \textbf{Linux 5.15.0} + ARMv7l + PlutoSDR Rev.A 标准 rootfs \\
烧录接口 & \texttt{/sys/class/fpga\_manager/fpga0/firmware}（write \texttt{.bin} filename） \\
\bottomrule
\end{tabularx}
\end{table}

\subsubsection{烧录验证（2026-05-14 上午）}

\begin{lstlisting}[language=bash]
$ echo ldsdr_2tr_emeye_safe.bin > /sys/class/fpga_manager/fpga0/firmware
WRITE_OK
$ dmesg | tail -2
[ ... ] fpga_manager fpga0: writing ldsdr_2tr_emeye_safe.bin to Xilinx Zynq FPGA Manager
[ ... ] fpga_manager fpga0: state=operating
\end{lstlisting}

烧录后稳定性验证清单：FPGA Manager \texttt{state = operating} \statusok；千兆 MAC 没断，SSH 连接保持 \statusok；板子稳定 30+ 分钟无 kernel oops \statusok；dmesg AD9361 错误零条 \statusok。

\subsubsection{AD9361 RF 链路验证（204 MHz EM Eye 频点）}

AD9361 配置到 EM Eye byte-clock 谐波频段并真实接收 RF：

\begin{table}[H]
\centering
\caption{AD9361 RF 链路参数}
\label{tab:ad9361-rf}
\small
\begin{tabularx}{\textwidth}{lXX}
\toprule
\rowcolor{tableheader}
\textbf{AD9361 参数} & \textbf{值} & \textbf{备注} \\
\midrule
ENSM mode & \texttt{fdd} & 工作模式 \\
默认 / 配置后 RX\_LO & 2.4 GHz / \textbf{204 MHz} & EM Eye byte-clock 谐波 \\
采样率 SR & \textbf{8 MSPS} & 与 §B.2 CIC 抽取一致 \\
RX Gain & 50 dB & manual gain control \\
\textbf{RSSI} & \textbf{93.75 dB} & 真实 RF 接收信号强度 \\
\bottomrule
\end{tabularx}
\end{table}

IIO buffer 采集 \textbf{16 KB 真实 IQ 数据}（signed 12-bit ADC），前 16 字节例：

\begin{lstlisting}[basicstyle=\ttfamily\footnotesize]
f0ff d5ff e4ff ffff ebff 1f00 ecff e1ff ...
\end{lstlisting}

非零，典型空气接收 baseline 噪声 + 信号（signed 12-bit little-endian：\texttt{0xfff0} \(\to -16\)，\texttt{0xffd5} \(\to -43\)）。RX path 真实导通。工具链 \texttt{iio\_attr} + sysfs 直接读写。

\textbf{提案意义}：Phase 1 W1 硬件 bring-up 风险被前置消除，LDSDR 进 lab 当天直接进入 Phase 1 Day 3（emeye\_accel IQ tap 接线），省 1-2 周。

\subsection{Phase 2 准备文档（提前完成）}
\label{sec:appB-5}

工程预案合计 1{,}096 行（451 + 251 + 394）。

\subsubsection{Phase 2 IQ Tap 方案}

\texttt{fpga\_accel/doc/phase2\_iq\_tap.md}（451 行）：BD 中 \texttt{axi\_ad9361} ADC AXI-Stream tap 到 \texttt{emeye\_accel\_top} 两方案对比（A 顶层手动接 / B BD 内 IP block 包装，推荐 B）；5 项待验证（时钟域、stream 位宽、tready 反压、空 idle、reset 序列）；Phase 2 W1 Day-by-Day（D1 tap \(\to\) D2 单通 \(\to\) D3 双通 \(\to\) D4 同步触发 \(\to\) D5 联调）。

\subsubsection{Phase 2 PS 用户态 TCP Forwarder}

\texttt{phase2\_tcp\_forwarder.md}（251 行）+ \texttt{.c}（394 行骨架）：AXI-DMA S2MM 寄存器全部标注 PG021 出处；UIO + mmap 用户态访问 AXI 寄存器（不写 kernel driver）；ping-pong buffer（2\(\times\)4 MB）避免 DMA 等 socket；TCP\_NODELAY + SO\_SNDBUF=4 MB；与 \texttt{simulation/network\_demo.py} 协议兼容（32-byte 帧头 + little-endian payload）。

Phase 2 W2 主机侧把 \texttt{simulation/network\_demo.py} 的 mock socket 接到 forwarder \(\to\) 管线闭环。

\subsection{工作量与产物清单}
\label{sec:appB-6}

\begin{table}[H]
\centering
\caption{产物清单}
\label{tab:artifacts}
\footnotesize
\begin{tabularx}{\textwidth}{lXcc}
\toprule
\rowcolor{tableheader}
\textbf{类别} & \textbf{文件} & \textbf{行数 / 大小} & \textbf{状态} \\
\midrule
\textbf{仿真} & \texttt{simulation/emeye\_simulation.py} & 501 行 & \statusok \\
              & \texttt{simulation/hdmi\_simulation.py} & 423 行 & \statusok \\
              & \texttt{simulation/network\_demo.py} & 405 行 & \statusok \\
              & \texttt{simulation/output/*.png} & 8 张图 & \statusok \\
\textbf{RTL 设计} & \texttt{fpga\_accel/rtl/magnitude\_jpl.v} & 117 行 & \statusok \\
              & \texttt{fpga\_accel/rtl/cic\_decimator.v} & 86 行 & \statusok \\
              & \texttt{fpga\_accel/rtl/frame\_sync.v} & 142 行 & \statusok \\
              & \texttt{fpga\_accel/rtl/emeye\_accel\_top.v} & 282 行 & \statusok \\
\textbf{Testbench} & \texttt{tb\_magnitude\_jpl.v} & 177 行 & \statusok\ 60/60 PASS \\
              & \texttt{tb\_cic\_decimator.v} & 222 行 & \statusok\ 11/11 PASS \\
              & \texttt{tb\_frame\_sync.v} & 176 行 & \statusok\ 5/5 PASS \\
              & \texttt{tb\_emeye\_accel.v} & 222 行 & \statusok\ PASS \\
\textbf{设计文档} & \texttt{fpga\_derivations.md} & 529 行 & \statusok \\
              & \texttt{phase2\_iq\_tap.md} & 451 行 & \statusok \\
              & \texttt{phase2\_tcp\_forwarder.md} & 251 行 & \statusok \\
              & \texttt{phase2\_tcp\_forwarder.c}（骨架） & 394 行 & \statuspartial \\
\textbf{Bitstream} & \texttt{ldsdr\_2tr\_emeye.bit} & 2.0 MB, MD5 \texttt{01a664a9...} & \statusok \\
              & \texttt{ldsdr\_2tr\_emeye\_safe.bin} & 2.0 MB, MD5 \texttt{b294f2a3...} & \statusok \\
              & \texttt{bit\_to\_bin.py} & 53 行 & \statusok \\
\textbf{总计} & 工程文件 \(\sim 20\) 份 & \textbf{\(\sim 6{,}000\) 行} & \\
\bottomrule
\end{tabularx}
\end{table}

状态分级：\statusok\ 已完成（代码存在、可运行/可综合、有可重现实测）；\statuspartial\ 部分完成（骨架/接口/方案已定，Phase 1-2 W1 内完成最后接线）；\statusplan\ Phase 1 计划（尚未动工）。

\subsection{这些工作支撑的提案 claim 验证}
\label{sec:appB-7}

\begin{table}[H]
\centering
\caption{提案 claim 支撑证据}
\label{tab:claim-evidence}
\footnotesize
\begin{tabularx}{\textwidth}{XXc}
\toprule
\rowcolor{tableheader}
\textbf{提案 claim} & \textbf{支撑证据} & \textbf{证据强度} \\
\midrule
Ch2 §2.3 CT 选型 + 多频段融合可行 & §B.3.2 LUT 13.75\%，Phase 3 86\% 余量 & \textbf{硬实测} \\
Ch2 §2.4 加速器纯整数 + 零 DSP & §B.2.1 magnitude\_jpl 仅移位+加 + §B.3.2 DSP = 0 & \textbf{硬实测} \\
Ch2 §2.5 跨时钟域风险已识别 + 缓解 & §B.3.3 WNS \(-2.985\) ns + Phase 1 W1 双方案 & \textbf{硬实测 + 明确计划} \\
Ch2 §2.6 JPL 误差在可接受区间 & §B.2.2 平均 4.29\% / 峰值 6.80\% & \textbf{硬实测} \\
Ch3 §3.1 EM Eye 论文方法已掌握 & §B.1.1 Python SSIM 0.98 & \textbf{硬实测} \\
Ch3 §3.4 平台可泛化到 HDMI TEMPEST & §B.1.2 HDMI SSIM 0.9907 / Corr 0.9976 & \textbf{硬实测} \\
Ch3 §3.5 LDSDR 已可控，Phase 1 风险低 & §B.4 真板 SSH + 烧录 + RSSI 93.75 dB & \textbf{硬实测} \\
Ch3 §3.6 AD9361 可调到 EM Eye 频段 & §B.4.3 RX\_LO 204 MHz + 16 KB IQ & \textbf{硬实测} \\
Ch4 §4.1 网络管线主机侧无 blocker & §B.1.3 network\_demo + §B.5.2 PS 骨架 & \textbf{硬实测 + 骨架就绪} \\
Ch4 §4.2 Phase 2 Day-by-Day 不卡壳 & §B.5.1 IQ tap 方案 + 5 项待验证清单 & \textbf{方案就绪} \\
\bottomrule
\end{tabularx}
\end{table}

\subsection{小结}
\label{sec:appB-8}

Phase 0 关键工程风险（论文方法理解、RTL 可综合性、目标器件资源容量、真板可控性、RF 链路在论文频点的可用性、Phase 2 主机管线协议）在提案提交之前已逐一前置消除。Phase 1 W1 第 1 天可直接进入 \texttt{emeye\_accel} 与 \texttt{axi\_ad9361} IQ tap 真板联调。

\section{风险与未决问题清单}
\label{sec:appC}

本附录诚实列出 \textbf{Phase 0 决策门之前尚未排除的工程风险}。每项均给出严重程度、缓解措施、与 Go/No-Go 关联。

\begin{table}[H]
\centering
\caption{Phase 0 前 10 条工程风险清单}
\label{tab:risk-list}
\footnotesize
\begin{tabularx}{\textwidth}{cXcXc}
\toprule
\rowcolor{tableheader}
\textbf{\#} & \textbf{风险描述} & \textbf{严重} & \textbf{缓解措施} & \textbf{决策点} \\
\midrule
1 & 电源线 100 MHz–1 GHz 传导衰减未知，实测前估 40–60 dB，超 60 dB 拖垮 SNR & 高 & Phase 0 W2 LISN + 频谱仪实测；LNA +28 dB 缓冲；退路近场探头补充 & Go/No-Go \#1 \\
2 & 铁氧体磁芯 CT 在 \(>500\) MHz 段磁导率下降，带宽不足 & 中 & 主选 Tekbox TBCP2-1000（实测 1 GHz）；辅路电容耦合补 \(>500\) MHz & Phase 1 W3 \\
3 & COTS 设备差异：Table II 12 台未必每台电源线泄漏强 & 中 & Phase 0 W3 优先复现 RPi V1；Phase 2 扫描 Wyze/Dafang 等 5 台 & Go/No-Go \#2 \\
4 & LDSDR/Pluto 千兆 Eth UDP/IIO 在持续 30 MSPS 下可能丢包 & 低 & 候选人 OFDM+LDPC BER=0 经验；千兆余量 25\(\times\)；退路 v1 走 8 MSPS，v2 板上加速 & Phase 1 W4 \\
5 & AD9363 \(\to\) AD9361 解锁固件偶尔失效 & 低 & 已验证 LDSDR rev2.1：RX\_LO 204 MHz · RSSI 93.75 dB · 16 KB IQ（详见附录 B） & N/A \\
6 & EMC Class B + IEC 61010 安全合规：市电耦合涉及隔离与浪涌防护 & 中 & Ch2 §2.4 GDT+MOV+共模扼流+TVS+LC HPF 五级链；Phase 1 W2 高压浪涌实测 & Phase 1 W2 \\
7 & 真 IQ tap \(\to\) \texttt{emeye\_accel} 跨时钟域（adc\_clk 100 MHz \(\to\) emeye\_clk 8 MHz） & 中 & 已生成 idle bitstream 验证综合通过；附录 B §B.5 已起草跨域方案 & Phase 1 W1 \\
8 & Vivado 实现 rx\_clk \(\to\) AXI 跨时钟域 WNS \(-2.985\) ns（原 ad9361 模板已有） & 低 & \texttt{set\_false\_path} 或 async FIFO 修复，与 \texttt{emeye\_accel} 无关 & Phase 1 W1 \\
9 & ADALM-Pluto / LDSDR 供应链：AD9361 国内供应偶尔受限 & 低 & 候选人 BYO LDSDR rev2.1，无需采购 & N/A \\
10 & PDF 最终交付：目标 25–30 页 + 完整插图 & 低 & 已经 1000+ 行 markdown，本周末追加 4 张系统框图 + 2 张测试图 & 2026-05-17 \\
\bottomrule
\end{tabularx}
\end{table}

\paragraph{Go/No-Go \#1 触发条件} Phase 0 W2 末 LISN 实测 100 MHz–1 GHz 耦合衰减 \(> 65\) dB \(\to\) 路线 A 申请更高增益 LNA（40–60 dB，单级 NF \(< 2\) dB）；路线 B 改聚焦近场 + 短电源线段（\(< 2\) m）。

\paragraph{Go/No-Go \#2 触发条件} Phase 0 W3 末 RPi V1 隔离变压器后端 SNR \(< 15\) dB \(\to\) 路线 A 换 DUT（Wyze Cam Pan / Dafang 1080P 等更强泄漏）；路线 B 增加被动 BPF SAW + 模拟下变频前置 LDSDR。

% ============================================================================
\section*{结语}
\addcontentsline{toc}{section}{结语}
% ============================================================================

\textbf{6 个月承诺}：6 个月全职投入，每周 GitHub commit，每月实验室汇报，phase 末 PI/Co-PI 评估。

\textbf{12 个月技术支持窗口期}：源码与实测数据开源至实验室仓库，Phase 3 由后继者完成期间提供 12 个月技术支持。

\vspace{2em}
\hrule
\vspace{0.5em}
\noindent\small\textit{致谢}：感谢 Long Y. \textit{et al.} (NDSS 2024) EM Eye 论文方法（"自相关 30 Hz 帧同步 + amplitude reshape"），本提案在不改变其核心算法的前提下将物理通道从空中辐射换为电源线传导。

\end{document}
